\section{Polynomial Approximation}

\noindent\textbf{Function evaluation.}
We have studied computable real numbers and floating-point arithmetic.
We now ask how to evaluate a continuous function
$f\colon\R\to\R$ at $x\in\mathbb F$ efficiently and to a prescribed
accuracy. Computability guarantees arbitrarily accurate approximations,
but does not by itself make their computation efficient.

For example, to evaluate $\exp(x)$ we may choose a Taylor degree $n$,
compute the rational Taylor approximation
\[
 1+x+\frac{x^2}{2!}+\cdots+\frac{x^n}{n!}
\]
exactly using big-integer arithmetic, and finally round to the target
floating-point format. This can be expensive: exact rational
arithmetic is slower than floating-point arithmetic, large $|x|$
may require large $n$, and near $x=0$ an unnecessarily large degree
may be suboptimal. The goal is therefore to choose an effective
approximation and evaluate it with inexpensive floating-point
operations. We will treat the approximation problem first,
assuming exact arithmetic; rounding during evaluation is a
separate issue.

\begin{definition}[Polynomial Approximation]
Let $a<b$ and let $C([a, b])$ denote the continuous real-valued
functions on $[a, b]$. In this section and the next, write
\[
 \begin{aligned}
 P_n&\coloneqq\left\{\sum_{k=0}^n c_kx^k\ \middle|\ c_k\in\R\right\},
       \qquad n\in\Nzero, \\
 P&\coloneqq\bigcup_{n=0}^{\infty}P_n.
 \end{aligned}
\]
For $f\in C([a, b])$ and $p\in P$, the \emph{uniform} and
\emph{$L^2$ approximation errors} are, respectively,
\[
 \|f-p\|_\infty
 =\max_{x\in[a, b]}|f(x)-p(x)|,
 \qquad
 \|f-p\|_2
 =\left(\int_a^b|f(x)-p(x)|^2\,\dd x\right)^{1/2}.
\]
We focus on uniform approximation: find $p\in P$ that makes
$\|f-p\|_\infty$ small, then evaluate $p$ on a computer.
\end{definition}

Polynomials are attractive because floating-point addition,
multiplication, and fused multiply-add are directly supported by
hardware. Rational functions may sometimes be even more efficient or
accurate; for instance, a term of the form
$1/(\pi/2-x)$ can capture the behavior of $\tan x$ near $\pi/2$.

The approximation problem raises three questions:
\begin{enumerate}[label=(\roman*)]
 \item \emph{Universal approximation:} can every
 $f\in C([a, b])$ be uniformly approximated arbitrarily well by polynomials?
 \item \emph{Best approximation:} for a fixed degree bound $n$,
 does a polynomial minimizing $\|f-p\|_\infty$ exist, and is it unique?
 \item \emph{Construction:} how can such a best polynomial be computed?
\end{enumerate}
The first question is answered here; the other two lead to
minmax polynomials.

\Needspace{12\baselineskip}
\begin{theorem}[Weierstrass Approximation Theorem]
For every $f\in C([a, b])$ and $\varepsilon>0$, there exists
$p\in P$ such that
\[
 \|f-p\|_\infty<\varepsilon.
\]
\end{theorem}
\begin{proof}
It suffices to work on $[0, 1]$. Indeed, for a given
$f\in C([a, b])$, define $g(t)=f(a+(b-a)t)$ on $[0, 1]$.
If $q$ approximates $g$ uniformly within $\varepsilon$,
then $p(x)=q((x-a)/(b-a))$ is a polynomial satisfying
$\|f-p\|_\infty<\varepsilon$.

Now let $f\in C([0, 1])$. For $n\ge1$ and $x\in[0, 1]$,
let $X_n$ satisfy $nX_n\sim\Bin(n, x)$.
Then
\[
 \Exp[X_n]=x,
 \qquad \Var(X_n)=\frac{x(1-x)}n
 \le\frac1{4n}.
\]
Consider the \emph{Bernstein polynomial}
\[
 \begin{aligned}
 f_n(x)
 &\coloneqq\Exp[f(X_n)]\\
 &=\sum_{k=0}^n\binom nk x^k(1-x)^{n-k}f(k/n).
 \end{aligned}
\]
Each $f_n$ is a polynomial of degree at most $n$.

Since $f$ is uniformly continuous on $[0, 1]$, choose
$\delta>0$ such that
$|s-t|<\delta$ implies $|f(s)-f(t)|<\varepsilon/2$.
Split the expectation according to
$|X_n-x|<\delta$ and $|X_n-x|\ge\delta$. Using
$|f(s)-f(t)|\le2\|f\|_\infty$ and Chebyshev's inequality,
we obtain, uniformly in $x$,
\[
 \begin{aligned}
 |f_n(x)-f(x)|
 &\le\Exp[|f(X_n)-f(x)|]\\
 &\le\frac\varepsilon2
    +2\|f\|_\infty\Prob(|X_n-x|\ge\delta)\\
 &\le\frac\varepsilon2
    +\frac{\|f\|_\infty}{2n\delta^2}.
 \end{aligned}
\]
Choose $n>\|f\|_\infty/(\varepsilon\delta^2)$.
Then $\|f-f_n\|_\infty<\varepsilon$, completing the proof.
\end{proof}
